\documentclass[12pt]{article}

% House style preamble
\input{.house-style/preamble.tex}
\usepackage{float}
\usepackage{placeins}

% Update PDF metadata
\hypersetup{
  pdftitle={Isocolon and comparative construal: A corpus test of rhetorical balance}
}

\title{Isocolon and comparative construal:\\ A corpus test of rhetorical balance}
\author{Brett Reynolds \orcidlink{0000-0003-0073-7195}%
\thanks{Contact: \href{mailto:brett.reynolds@humber.ca}{brett.reynolds@humber.ca}}\\
Humber Polytechnic \& University of Toronto}
\date{\today}

\begin{document}
\maketitle

\begin{abstract}
Rhetorical theory often treats \term{isocolon} as more than ornamental
balance. A balanced pair of cola can make two units feel more explicitly
comparable. A reader-side causal claim would need experimental evidence,
but a production-side textual claim can be tested in corpus data. This
paper tests whether adjacent discourse-relation pairs in GUM/eRST are
more isocolonic when they construe two units as comparable. The outcome
is a continuous score combining length balance, syntactic parallelism,
and lexical parallelism, rather than a binary figure label. In 17,456
adjacent relation pairs from 299 documents across 24 genres, a broad
comparison and co-ordination family shows a clear adjusted increase in
isocolonicity, especially for lists, disjunctions, and adversative
contrast. A narrow adversative family is smaller, antithesis alone is
near zero, and concession is negative. The result supports a restrained
textual claim: formal balance is part of the profile of some comparative
construals, but it is not a general property of every relation that
could be called contrastive.
\end{abstract}

\section{Introduction}

Rhetorical theory has long treated figures as patterned forms with
argumentative force, not as detachable ornaments. In that tradition,
\term{isocolon} is interesting because it joins two kinds of evidence
that are easy to confuse: formal balance between adjacent units and a
reader's sense that the units invite comparison. The motivating claim is
simple. When writers place two clauses or phrases in balanced form, they
may make the comparison between them more explicit, pointed, or forceful.

This paper asks how much of that claim can be tested without a reader
experiment. The answer depends on the explanatory level. At the
reader-response level, the claim is causal: isocolon changes what readers
notice or how forcefully they construe a relation. Corpus evidence alone
can't show that. At the textual level, the claim is distributional:
writers may preferentially use balanced form where they construe two
units as comparable. That claim is testable.

I therefore test a narrower hypothesis. When a relation places two
adjacent units in a comparison-like configuration, the pair should be more
isocolonic than pairs expressing cause, condition, elaboration, sequence,
or other non-comparative relations. The study uses the GUM corpus and
eRST annotations because they provide genre-diverse texts,
discourse-relation labels, and signal annotation
\citep{zeldes2017GUMCorpus,zeldesEtAl2025eRST}.
The current analysis is preliminary because the RST Signalling Corpus
data are still pending, but the design already gives a useful first test.

The results are mixed in the informative sense. Broadly comparative and
co-ordinative relations are more formally balanced than matched
non-target relations. The effect is clearest for \mention{joint-list},
\mention{joint-disjunction}, and \mention{adversative-contrast}. It is
much weaker for a narrow adversative set, near zero for
\mention{adversative-antithesis} alone, and negative for
\mention{adversative-concession}. The pattern refines the rhetorical
claim: isocolonicity tracks some acts of comparative construal better
than it tracks the inherited relation name
\term{antithesis}.

\section{Rhetorical figures as form-function pairings}

The theoretical starting point is Fahnestock's account of rhetorical
figures as form-function pairings. In her work on scientific argument and
rhetorical style, figures are not decorative additions to an argument;
they are recurrent formal arrangements that can epitomize lines of
reasoning \citep{fahnestock1999RhetoricalFiguresScience,fahnestock2011RhetoricalStyle}.
That view matters here because it makes a corpus prediction. If a figure
is tied to an argumentative operation, then the figure should be
distributed unevenly across the discourse contexts where that operation
is being performed.

Fahnestock's discussion of the figures of argument is especially direct
for this project. \term{Isocolon} concerns approximate equality in the
length of adjacent cola, while \term{parison} concerns similarity of
grammatical structure. The two often co-occur, but they are not the same
property. A pair can be close in length without being syntactically
parallel, and it can be syntactically parallel without being closely
balanced in length. The classical tradition links both to comparison and
contrast because similar form can make two units cohere as an
argumentative set \citep{fahnestock2011RhetoricalStyle}.

The same point motivates computational work on rhetorical figures.
\textcite{harrisDimarco2017RhetoricalFiguresComputation} argue that
figural form is often iconically related to discourse function, so
surface patterning can provide evidence about pragmatic and argumentative
actions. \textcite{harrisEtAl2018AnnotationSchemeRhetoricalFigures}
develop an annotation scheme for figures in part because figure
identification requires attention to both formal regularity and
rhetorical role. \textcite{lawrenceEtAl2017HarnessingRhetoricalFigures}
make the parallel point for argument mining: rhetorical figures can
identify regions where argumentative work is being done, but only if
formal detection is connected to the content and relation being advanced.

This literature supports the central assumption of the present study.
The relevant object is the co-occurrence of surface balance with a
discourse relation that makes two units available for comparison. The
operational question is whether a graded measure of balance is higher in
the relations where rhetorical theory predicts that balance should be
useful.

\section{From candidate form to rhetorical function}

The main methodological risk is false positives. Repetition and
parallelism are common in ordinary text, and only some instances have the
rhetorical force that figure theory cares about. Work on automatic
chiasmus detection makes the problem vivid. \textcite{dubremetzNivre2018FigureDetection}
show that trivial repetition patterns are easy to find, but repetitions
with rhetorical effect are much harder to separate from accidental or
irrelevant patterns. They therefore treat figure detection as a ranking
problem rather than a clean binary classification problem.

That lesson carries over to isocolon. Exact repetition, list formatting,
dialogue formulas, quoted slogans, and segmentation artifacts can all
look formally balanced. Some are rhetorically relevant, but many are not
instances of the target phenomenon. Recent survey work makes the broader
point: rhetorical figure detection remains limited by small datasets,
inconsistent definitions, and over-reliance on rule-based surface cues
\citep{kuehnEtAl2024LesserKnownRhetoricalFigures}. A credible corpus
test should therefore use surface scores cautiously, validate them
against human judgment, and avoid treating high balance as equivalent to
a gold figure label.

The analysis below follows that strategy. The isocolonicity score is a
continuous predictor-like outcome: it ranks adjacent relation pairs by
how much balance they display. The qualitative audit then asks whether
high-scoring examples are plausibly isocolonic after excluding exact
repetition, dialogue formulas, and other false-positive types. This is a
more modest target than automatic figure recognition, but it is the
right target for a first rhetorical-function test.

\section{Corpus and measurement}

The corpus is the English GUM/eRST dataset. GUM is a genre-diverse
corpus with richly annotated English texts
\citep{zeldes2017GUMCorpus}. eRST extends RST-style discourse annotation
with graph structures and explicit signal annotation across spoken and
written genres \citep{zeldesEtAl2025eRST}. This matters because the
project needs discourse relations rather than only adjacent sentences. It
also needs enough genre variation to separate a rhetorical effect from a
register-specific habit.

The unit of analysis is an adjacent discourse-relation pair. The current
dataset contains 17,456 such pairs from 299 documents across 24 genres.
For each pair, I computed a composite isocolonicity score from three
families of surface evidence. The first is length balance, using token
length and related measures of how similar the two units are in size. The
second is syntactic parallelism, based on how similar their dependency
and part-of-speech profiles are. The third is lexical parallelism, based
on repeated or closely corresponding lexical material. The composite
should be read as a graded measure of formal balance, rather than as a
claim that the pair is an isocolon.

The score is deliberately simple. For any two units with positive length,
a balance component is \(1 - |n_1 - n_2|/\max(n_1,n_2)\). The length
score averages this quantity over word tokens, characters, and estimated
syllables. The syntax score averages normalized edit similarity over
UPOS, XPOS, and dependency-relation sequences. The lexical score averages
lemma-sequence edit similarity and Jaccard overlap over content lemmas.
The composite is
\[
0.40 \times \text{length} +
0.45 \times \text{syntax} +
0.15 \times \text{lexical}.
\]
These weights give syntax slightly more influence than length, with
lexical overlap included as a weaker cue.

The target relations were defined in two ways. The broad target family
includes \mention{adversative-antithesis},
\mention{adversative-contrast}, \mention{joint-disjunction}, and
\mention{joint-list}. This family tests whether formal balance rises in
relations that place units in a comparable or co-ordinated frame. The
narrow target family includes only \mention{adversative-antithesis} and
\mention{adversative-contrast}. Individual-relation analyses then separate
the broad effect into its component labels.

The main comparisons adjust for the most obvious structural confounds.
Where possible, pairs are compared within documents. The models also
include controls for explicit signalling, sentence adjacency,
punctuation, mean length, and maximum length. A genre-varying analysis
then checks whether the effect is concentrated in one or two genres or
survives partial pooling across genres.

The adjusted estimates are ordinary least-squares target coefficients
computed by residualizing both the target indicator and the outcome
against the controls. The controls are an intercept, explicit relation
status, same-sentence adjacency, colon punctuation, semicolon
punctuation, standardized log mean length, standardized log maximum
length, and document fixed effects. The genre-varying analysis first fits
the same adjusted effect within each genre, then uses a normal
empirical-Bayes model with genre effects distributed around a shared mean
and between-genre standard deviation.

\section{Validation}

The score behaves as a useful ranking measure, not as a binary detector.
Against the small number of eRST \mention{syn-prl} signal labels in the
dataset, the composite score has an AUC of 0.731. The syntax subscore is
slightly higher at 0.754, the lexical subscore is 0.719, and the length
subscore is 0.646 (Table~\ref{tab:score-validation}). At the 95th percentile,
precision is only 3.1\%, but
recall is 37.0\%, a lift of 7.4 over prevalence. At the 97.5th
percentile, precision rises to 5.5\%, recall is 32.9\%, and lift is
13.1. Those numbers are poor for automatic labelling but useful for
ranking rare candidates for audit.

\begin{table}[H]
\centering
\small
\begin{tabular}{lrrr}
\toprule
Score & AUC & Gold mean & Non-gold mean\\
\midrule
Composite & 0.731 & 0.502 & 0.327\\
Length & 0.646 & 0.683 & 0.560\\
Syntax & 0.754 & 0.441 & 0.209\\
Lexical & 0.719 & 0.202 & 0.057\\
\bottomrule
\end{tabular}
\caption{Score validation against the 73 eRST \mention{syn-prl} signal labels
in 17,456 adjacent relation pairs.}
\label{tab:score-validation}
\end{table}

The qualitative audit is therefore central. Eighty candidate pairs were
human-coded, with independent model-assisted reviews used only to surface
borderline cases. The combined human coding yielded 33 yes, 43 no, and 4
uncertain judgments. High-scoring broad-target candidates performed well:
the first batch had 8 yes judgments out of 10 high broad-but-not-narrow
items, and the second had 6 yes judgments out of 6 high broad-target
items. By contrast, the restatement/repetition controls had 0 yes
judgments out of 9. This contrast is important because exact repetition
is the strongest surface correlate of formal balance, while also being a
separate false-positive class for the purposes of this study.

The audit also clarified the boundary of the category. Bare identity,
apposition, dialogue repair, and exact repetition can be rhetorically
relevant without being isocolon. Conversely, weak length balance can
still matter when the two cola share a frame strongly enough to invite
slot substitution. The coding decisions therefore treat isocolon as a
graded rhetorical judgment with false-positive classes rather than a
mechanically recoverable surface pattern.

\FloatBarrier

\section{Preliminary results}

The broad target family shows a clear adjusted increase in composite
isocolonicity. Across all adjacent pairs, the adjusted effect is 0.0415
with a confidence interval from 0.0375 to 0.0455. The subscales show that
the effect is driven most strongly by syntax: the adjusted syntax effect
is 0.0770, compared with 0.0244 for lexical parallelism and 0.0080 for
length balance. That pattern is theoretically useful. The corpus effect
reflects shared formal frames more than similar size. Table~\ref{tab:relation-effects}
summarizes the main adjusted effects and the restatement positive control.

The narrow adversative family is smaller. For
\mention{adversative-antithesis} plus \mention{adversative-contrast}, the
adjusted composite effect is 0.0130, with a confidence interval from
0.0058 to 0.0202. The individual labels explain the difference.
\mention{joint-disjunction} has an adjusted effect of 0.0482,
\mention{joint-list} has an adjusted effect of 0.0439, and
\mention{adversative-contrast} has an adjusted effect of 0.0241.
\mention{adversative-antithesis}, despite its rhetorical name, is near
zero at 0.0039, with an interval that crosses zero.
\mention{adversative-concession} is negative at -0.0131.

The relation atlas also identifies an important positive control.
\mention{restatement-repetition} has the strongest adjusted composite
effect, 0.1902. That is exactly what a surface score should find: repeated
material is highly balanced. But the human audit shows why this can't be
the target effect. Repetition can signal clarification, identity,
repair, emphasis, or other rhetorical work; exact repetition remains
distinct from balanced comparison. The main analysis therefore
treats restatement/repetition as a control and diagnostic, not as
evidence for the comparison hypothesis.

Genre-varying estimates point in the same direction. The broad target
has a partially pooled mean effect of 0.0390 across 20 genres. The
\mention{joint-list} estimate is 0.0398 across 17 genres, and
\mention{adversative-contrast} is 0.0291 across the genres where there
is enough evidence to estimate it. The narrow adversative estimate is
0.0118. These estimates are not identical across genres, but the broad
effect isn't confined to a single genre.

\begin{table}[H]
\centering
\small
\begin{tabular}{lrrr}
\toprule
Target & Target n & Target mean & Adjusted effect [95\% CI]\\
\midrule
Broad family & 2,462 & 0.401 & 0.041 [0.038, 0.045]\\
Narrow adversative & 605 & 0.382 & 0.013 [0.006, 0.020]\\
\mention{joint-list} & 1,690 & 0.407 & 0.044 [0.039, 0.049]\\
\mention{joint-disjunction} & 167 & 0.409 & 0.048 [0.035, 0.062]\\
\mention{adversative-contrast} & 262 & 0.408 & 0.024 [0.013, 0.035]\\
\mention{adversative-antithesis} & 343 & 0.363 & 0.004 [-0.006, 0.013]\\
\mention{adversative-concession} & 679 & 0.342 & -0.013 [-0.020, -0.006]\\
\mention{restatement-repetition} & 321 & 0.580 & 0.190 [0.181, 0.200]\\
\bottomrule
\end{tabular}
\caption{Selected adjusted effects on composite isocolonicity. Estimates use
document fixed effects and controls for explicitness, sentence adjacency,
punctuation, and length. The restatement row is a positive control rather than
a comparison target.}
\label{tab:relation-effects}
\end{table}

\FloatBarrier

\section{Interpretation}

The findings support a restricted version of the rhetorical-function
claim. Formal balance is elevated where writers construe adjacent units
as members of a comparable set, especially in lists, disjunctions, and
ordinary adversative contrast. The evidence is weaker for labels that
look rhetorically relevant in the taxonomy but don't necessarily require
balanced presentation in the text. In particular,
\mention{adversative-antithesis} as an eRST relation label only partly
overlaps with antithesis in the classical theory of figures.

This distinction matters. If the analysis had simply shown that
\mention{adversative-antithesis} was high, the result would have been
easy to summarize but theoretically shallow. The actual pattern is more
useful. Isocolonicity appears to be associated with a broader act of
placing units into a comparable frame. That frame can be contrastive, but
it can also be co-ordinative, disjunctive, or list-like. Conversely, some
adversative relations are not especially isocolonic because their
function is concession, qualification, or discourse management rather
than balanced comparison.

The result also clarifies the scope of the corpus study. The evidence is
textual and distributional, rather than psycholinguistic: it concerns
where balanced form occurs, not how readers experience those pairs. Nor
does a high score by itself establish a rhetorical figure. The result
shows instead that a graded measure of balance is systematically higher
in particular discourse-relation contexts, and that human audit supports
the rhetorical relevance of many high-scoring broad-target cases.

The next robustness step is to add the RST Signalling Corpus once the
data are available. Work on discourse signalling has shown that coherence
relations are often signalled by more than discourse markers, including
lexical, syntactic, graphical, semantic, and genre-based cues. That point
is developed in the original signalling annotation work and in later
studies of RST signals and multiple signals
\citep{taboadaDas2013AnnotationUponAnnotation,dasTaboada2018RSTSignallingCorpus,dasTaboada2018SignallingBeyondMarkers,dasTaboada2019MultipleSignals}.
This literature gives the project a stronger comparative baseline.
If formal balance behaves like a discourse signal, then it should be
testable alongside other signal types, not only as an isolated figure
score.

\section{Conclusion}

A corpus study can't prove that isocolon makes readers perceive
comparison as more forceful. It can test a cleaner textual claim:
writers use formal balance unevenly across discourse relations, and the
relations most associated with comparable construals should show more of
it. The present results support that claim for a broad family of
comparison-like and co-ordinative relations, with the strongest evidence
coming from lists, disjunctions, and adversative contrast.

The result gives only limited support to the inherited figure label
\term{antithesis}. The narrow adversative effect is small, and antithesis
as a discourse-relation label is close to zero. That negative finding
helps sharpen the theory. Isocolon is a formal resource that writers can
use when two units are presented as comparable members of an
argumentative or expository set.

\section*{Acknowledgements}

This paper was drafted with the assistance of large language models. All
content has been reviewed and revised by the author, who takes full
responsibility for the final text.

\newpage
\printbibliography

\end{document}
